%%%%%%%%%%%%%%%%%%%%%%%%%%%%%%%%%%%%%%%%%
\chapter{Advanced Architectures}\label{advanced.architectures.ch}
%%%%%%%%%%%%%%%%%%%%%%%%%%%%%%%%%%%%%%%%%
So far, we have studied multi-layer perceptrons (MLPs), which are fully connected (dense) feedforward neural networks.  We have also seen that these networks perform well for a variety of pattern recognition problems that can be viewed as function approximation tasks. This is consistent with their property as universal approximators. We also studied their role in scientific computing through Physics-Informed Neural Networks (PINNs).


In this chapter, we focus on some special architectures that are designed for specific types of data or applications, like convolutional neural networks (CNNs), which are particularly suitable for image recognition, and recurrent neural networks (RNNs), which are suitable for language and sequence modelling.  We conclude the chapter with a discussion of a recent development called \hl{Transformers}, which is an attention-based architecture that has revolutionized sequence modelling. Transformers provide a unified framework for learning from text, images, and beyond.

In \sect{convolution.neural.network.sec}, we introduce convolutional neural networks. We first define the one-dimensional convolutional layer, together with stride, padding and pooling. We then extend these ideas to images and illustrate them with blurring and edge-detection kernels. In \sect{transformers.sec}, we turn to sequence data. We first recall recurrent neural networks briefly and point out their limitations. We then introduce self-attention, positional encoding and causal attention. Finally, we assemble these components into the transformer block and use it to build a model for sequence prediction.
%%%%%%%%%%%%%%%%%%%%%%%%%%%%%%%%%%%%%%%%%
%%%%%%%%%%%%%%%%%%%%%%%%%%%%%%%%%%%%%%%%%
\section{Convolutional Neural Networks}\label{convolution.neural.network.sec}
%%%%%%%%%%%%%%%%%%%%%%%%%%%%%%%%%%%%%%%%%
The dense structure of MLPs treats each input neuron independently and ignores the spatial structure of the input data. For example, in image data split into pixels, nearby pixels are highly correlated, and certain local features like edges or corners can appear at different locations in the image. The dense structure also leads to a very large number of parameters, which grows in proportion to the number of pixels. Convolutional Neural Networks (CNNs) overcome these limitations by incorporating the local spatial structure of data through \hl{local connectivity} and \hl{weight sharing}.

In local connectivity, each neuron is connected only to a small spatial region (patch) of the input in order to capture local patterns.  Through weight sharing, the same set of weights, referred to as a filter or kernel,  is used across different spatial locations. Weight sharing reduces the number of parameters. It also makes the layer translation equivariant, which means that a shift of the input produces the same shift of the output.  Pooling layers are often added to reduce spatial dimensions and to retain the most important features.  In this section, we briefly outline CNNs.
%%%%%%%%%%%%%%%%%%%%%%%%%%%%%%%%%%%%%%%%%
\subsection{One-Dimensional Convolutional Layer}
%%%%%%%%%%%%%%%%%%%%%%%%%%%%%%%%%%%%%%%%%
In many applications, such as time-series analysis and speech recognition, the input is indexed by a single variable, namely, time or position in the sequence.  At each index, the input may contain one or several values (channels).  For such data, \hl{one-dimensional convolutional neural networks} (1D CNNs) provide an efficient architecture that captures local dependencies along the temporal or sequential axis.

Let $f, g : \mathbb{R} \to \mathbb{R}$. 
The \hl{continuous one-dimensional convolution} of $f$ and $g$, denoted by $(f * g)$, is defined as
$$
(f * g)(x) = \int_{\mathbb{R}} f(\xi)\, g(x - \xi)\, d\xi,
$$
whenever the integral exists.

Analogously, in the discrete case where the signals are sequences, say, $f,g:\mathbb{Z}\rightarrow \R$ given by
$$
f = \{f_i\}_{i \in \mathbb{Z}}, 
\qquad 
g = \{g_i\}_{i \in \mathbb{Z}},
$$
the \hl{discrete convolution} is defined by
$$
(f * g)_i = \sum_{j \in \mathbb{Z}} f_j \, g_{i - j},
$$
whenever the series converges.
%%%%%%%%%%%%%%%%%%%%%%%%%%%%%%%%%%%%%%%%%
\begin{remark}[Finitely Supported Signals]

In practice, we have signals of finite length.  In this case, we extend the signal by zero. Such a signal is called a signal with \hl{finite support}.  For instance, if we have 
$$\bfm{g}=[g_1, g_2, \ldots, g_{N}],$$ 
for some $N>0$, then it should be understood as extending it by zero on either side and considering the signal of infinite length $g=[\ldots, 0, \bfm{g}, 0, \ldots]$.
\end{remark}
%%%%%%%%%%%%%%%%%%%%%%%%%%%%%%%%%%%%%%%%%
In CNNs, it is customary to use the \hl{cross-correlation} operation (instead of strict convolution), given by
$$
(\bfm{w} \star \bfm{x})_i = \sum_{j = 1}^{r}  w_j\, x_{i + j - 1}, ~~
i = 1, 2, \ldots, N - r + 1,
$$
where $\bfm{w} = [w_1, w_2, \ldots, w_{r}]$ is the \hl{kernel} or \hl{filter} of size $r$ and $\bfm{x}=[x_1, x_2, \ldots, x_{N}]$ is a finite input signal with $0 < r \le N$. In the deep learning literature, the cross-correlation is also referred to as convolution. We follow this convention in the CNN context below.
%%%%%%%%%%%%%%%%%%%%%%%%%%%%%%%%%%%%%%%%%

%%%%%%%%%%%%%%%%%%%%%%%%%%%%%%%%%%%%%%%%%
\begin{example}[Moving Average as a Convolution Operation]
Let 
$
\bfm{x} = [x_1,\, x_2,\, \ldots,\, x_N]
$
be a discrete-time signal (or sequence).  
A \hl{moving average} with window size $r$ computes each output value as the average of $r$ consecutive inputs, given by
$$
y_i = \frac{1}{r}\sum_{j=1}^{r} x_{i+j-1},
\quad i = 1, 2, \ldots, N - r + 1.
$$

This can be viewed as a \hl{cross-correlation} (or convolution without kernel flipping) between the input $\bfm{x}$ and the kernel
$$
\bfm{w} = \frac{1}{r}[1,\, 1,\, \ldots,\, 1],
$$
that is,
$$
y_i = \sum_{j=1}^{r} w_j\, x_{i+j-1}.
$$
Thus, the moving average is a special case of convolution where all kernel weights are equal and sum to one.
\end{example}
%%%%%%%%%%%%%%%%%%%%%%%%%%%%%%%%%%%%%%%%%
The one-dimensional moving average filter discussed above can be viewed as a \hl{low-pass filter}, since it smooths rapid variations in the input signal. In two dimensions, the same idea extends naturally to image processing, where a smoothing kernel corresponds to \hl{blurring} an image. 

In contrast, \hl{edge detection} in images corresponds to emphasizing changes (differences) in the signal, i.e., high-frequency components.
In one dimension, this can be seen directly through convolution with a \hl{difference kernel} or \hl{high-pass filter}.
%%%%%%%%%%%%%%%%%%%%%%%%%%%%%%%%%%%%%%%%%
\begin{example}[One-Dimensional Edge Detection]

Consider a one-dimensional input signal
$
\bfm{x} = [x_1,\, x_2,\, \ldots,\, x_6]
$
and a \hl{difference kernel} (or \hl{edge detection kernel})
$$
\bfm{w} = [-1, 1]
$$
The convolution (or, in CNNs, the cross-correlation) of $\bfm{x}$ and $\bfm{w}$ produces the output
$$
y_i = (-1)\,x_i + (1)\,x_{i+1} = x_{i+1} - x_i, 
\qquad i = 1, 2, \ldots, 5.
$$
Hence,
$$
\bfm{y} = [x_2 - x_1,\, x_3 - x_2,\, x_4 - x_3,\, x_5 - x_4,\, x_6 - x_5].
$$
This output measures the \hl{change} (or discrete gradient) between successive input values.
Regions where $\bfm{x}$ is approximately constant produce small responses, while sharp transitions yield large responses.  Such transitions correspond to \hl{edges} in one-dimensional data.

In two dimensions, similar difference kernels (e.g., Sobel or Prewitt filters) are used to detect edges in images.
\end{example}
%%%%%%%%%%%%%%%%%%%%%%%%%%%%%%%%%%%%%%%%%
\subsubsection{Convolutional Layer}
%%%%%%%%%%%%%%%%%%%%%%%%%%%%%%%%%%%%%%%%%
Let the input to a one-dimensional convolutional layer be a sequence (or signal)
$$
\bfm{x} = [x_1, x_2, \ldots, x_N] \in \mathbb{R}^N,
$$
and let the convolutional kernel (or filter) be
$$
\bfm{w} = [w_1, w_2, \ldots, w_r] \in \mathbb{R}^r,
$$
where the integer $1\le r\le N$ is called the \hl{kernel size} or \hl{receptive field length}. 
%%%%%%%%%%%%%%%%%%%%%%%%%%%%%%%%%%%%%%%%%
\begin{definition}[One-Dimensional Convolutional Layer]
For a given integer $s \ge 1$, a \hl{one-dimensional convolutional layer} with activation function $\actf$, convolution kernel $\bfm{w}$ and bias $b$ produces the output
\[
a_i = \actf\!\Bigg(\sum_{j=1}^{r} w_j\, x_{(i-1)s + j} +\,b\Bigg),
\quad i = 1, 2, \ldots, n_{\text{out}},
\]
where
\[
n_{\text{out}} = \Big\lfloor \frac{N - r}{s} \Big\rfloor + 1.
\]
The number $s$ is called the \hl{stride}, which specifies the step size with which the kernel moves along the input.
\end{definition}
%%%%%%%%%%%%%%%%%%%%%%%%%%%%%%%%%%%%%%%%%

%%%%%%%%%%%%%%%%%%%%%%%%%%%%%%%%%%%%%%%%%
The following example illustrates a one-dimensional convolutional layer with stride $s=1$. Such a layer can be expressed as a layer of a feedforward neural network with two constraints on its weights. Many weights are set to zero and are kept at zero during training. The remaining weights are shared by all neurons of the layer. A general result is stated in \prob{CNNs1FNN.pro}.
%%%%%%%%%%%%%%%%%%%%%%%%%%%%%%%%%%%%%%%%%
\begin{example}[Convolutional Layer with Stride $s=1$]
Consider an input sequence
$$
\bfm{x} = [x_1,\, x_2,\, x_3,\, x_4,\, x_5,\, x_6],
$$
and a kernel (filter)
$$
\bfm{w} = [w_1,\, w_2].
$$
Here, the input dimension is $N=6$, and the kernel size is $r=2$. Let $b$ be the bias.

Let us take stride $s = 1$. Then,
$$
n_{\text{out}} = \Big\lfloor \frac{N - r}{s} \Big\rfloor + 1
= \Big\lfloor \frac{6 - 2}{1} \Big\rfloor + 1 = 5.
$$
Therefore, from the convolutional layer definition, we have
$$
a_i = \actf\!\Bigg(\sum_{j=1}^{2} w_j\, x_{(i-1) + j}+\,b\Bigg), 
\quad i = 1, 2, \ldots, 5.
$$
The layer output activations are given by
$$
\begin{aligned}
a_1 &= \actf(w_1 x_1 + w_2 x_2 +\,b),\\
a_2 &= \actf(w_1 x_2 + w_2 x_3 +\,b),\\
a_3 &= \actf(w_1 x_3 + w_2 x_4 +\,b),\\
a_4 &= \actf(w_1 x_4 + w_2 x_5 +\,b),\\
a_5 &= \actf(w_1 x_5 + w_2 x_6 +\,b).
\end{aligned}
$$
Hence, the output feature map (activation vector) of the convolutional layer is
$$
\bfm{a} = [a_1,\, a_2,\, a_3,\, a_4,\, a_5].
$$
Writing $\bfm{x}$ and $\bfm{a}$ as column vectors, we get
$$
\bfm{a} = \vactf\big(W\bfm{x} + b\,\bfm{1}\big),
\qquad
W =
\begin{bmatrix}
w_1 & w_2 & 0 & 0 & 0 & 0\\
0 & w_1 & w_2 & 0 & 0 & 0\\
0 & 0 & w_1 & w_2 & 0 & 0\\
0 & 0 & 0 & w_1 & w_2 & 0\\
0 & 0 & 0 & 0 & w_1 & w_2
\end{bmatrix}\in\R^{5\times 6},
$$
where $\bfm{1}\in\R^5$ is the vector of ones. This is the matrix form of a layer of a feedforward network. Most entries of $W$ are zero, and the nonzero entries are shared by all rows.
\figu{conv1d.stride1.fig} shows this layer as a network diagram.
\end{example}
%%%%%%%%%%%%%%%%%%%%%%%%%%%%%%%%%%%%%%%%%
\begin{figure}[ht]
\begin{center}
\input{Figures/Ch08CNN1DLayer.png}
\end{center}
\caption{The one-dimensional convolutional layer with kernel $\bfm{w}=[w_1, w_2]$ and stride $s=1$, drawn as a layer of a feedforward network. Each neuron $a_i$ is connected only to the inputs $x_i$ and $x_{i+1}$. All neurons share the weights $w_1$, $w_2$ and the bias $b$. All missing connections have weight zero.}\label{conv1d.stride1.fig}
\end{figure}
%%%%%%%%%%%%%%%%%%%%%%%%%%%%%%%%%%%%%%%%%
The convolutional layer replaces the affine transformation $\bfm{w}\cdot\bfm{x}+\,b$ of a neuron in an MLP by the local operation $\bfm{w} \star \bfm{x} + b$. Each output activation (feature map entry) $a_i$ can still be viewed as a neuron. However, this neuron is connected only to the $r$ inputs $x_{(i-1)s+1}, \ldots, x_{(i-1)s+r}$, which form its local receptive field. Moreover, all these neurons share the same weights $\bfm{w}$ and the same bias $b$. Nevertheless, the convolution operation can be equivalently expressed as a matrix-vector multiplication,  where the kernel $\bfm{w}$ is expanded into a structured sparse matrix (a Toeplitz matrix when $s=1$).  Thus, a convolutional layer can be viewed as a layer of a feedforward network with a constrained weight matrix.
%%%%%%%%%%%%%%%%%%%%%%%%%%%%%%%%%%%%%%%%%

%%%%%%%%%%%%%%%%%%%%%%%%%%%%%%%%%%%%%%%%%
\begin{problem}[Convolution Layer as a Fully Connected Layer]
\label{CNNs1FNN.pro}
Show that a one-dimensional convolutional layer can be expressed as a fully connected (dense) layer with a specially structured weight (sparse Toeplitz) matrix, when the stride is $s=1$.
\end{problem}
%%%%%%%%%%%%%%%%%%%%%%%%%%%%%%%%%%%%%%%%%

%%%%%%%%%%%%%%%%%%%%%%%%%%%%%%%%%%%%%%%%%
\begin{note}
For a stride $s>1$, row $i$ of the weight matrix starts at column $(i-1)s+1$. The matrix is still sparse and has shared weights. However, it is no longer a Toeplitz matrix.
\end{note}
%%%%%%%%%%%%%%%%%%%%%%%%%%%%%%%%%%%%%%%%%
\begin{problem}
Consider an input vector $\bfm{x} = [x_1, x_2, x_3, x_4, x_5, x_6]$, a kernel (filter)
$\bfm{w} = [w_1, w_2]$, and a bias term $b$.
Using the convolutional layer with stride $s = 2$ and activation function $\actf$,
compute the output activations.
\end{problem}
%%%%%%%%%%%%%%%%%%%%%%%%%%%%%%%%%%%%%%%%%

%%%%%%%%%%%%%%%%%%%%%%%%%%%%%%%%%%%%%%%%%
\begin{remark}[Padding] 
The number of zeros $p$ added at each end of the input sequence before performing convolution is called \hl{padding}.
Padding allows the kernel to be applied at the two ends of the input, so that the boundary entries contribute to more outputs. 
 
Given an input of length $N$, if we pad with $p$ zeros on each side, the effective input length becomes $N + 2p$.  Thus, for a kernel size $r$, a stride $s$, and padding $p$, 
the convolutional layer is given by
$$
a_i = \actf\!\Bigg(\sum_{j=1}^{r} w_j\, x_{(i-1)s + j - p}+\,b\Bigg),
\quad i = 1, 2, \ldots, n_{\text{out}},
$$
where
$$
n_{\text{out}} = \Bigg\lfloor \frac{N + 2p - r}{s} \Bigg\rfloor + 1.
$$
Here, we use the convention $x_k = 0$ for $k \le 0$ and for $k > N$.

With stride $s=1$ and $p=(r-1)/2$ for an odd kernel size $r$, the output length equals the input length $N$.
\end{remark}
%%%%%%%%%%%%%%%%%%%%%%%%%%%%%%%%%%%%%%%%%
\begin{problem}
Consider an input vector $\bfm{x} = [x_1, x_2, x_3, x_4, x_5, x_6]$, 
a kernel $\bfm{w} = [w_1, w_2]$, and a bias term $b$.
Using the convolutional layer with stride $s = 1$, 
activation function $\actf$, and padding of one zero on each side 
(\emph{i.e.}, $p=1$), compute the output activations.
\end{problem}
%%%%%%%%%%%%%%%%%%%%%%%%%%%%%%%%%%%%%%%%%

%%%%%%%%%%%%%%%%%%%%%%%%%%%%%%%%%%%%%%%%%
%%%%%%%%%%%%%%%%%%%%%%%%%%%%%%%%%%%%%%%%%
\begin{remark}[Pooling Layer]

After a convolutional layer produces its output (called the \hl{feature map}), it is often followed by a \hl{pooling layer}. 
Given an input feature map
$$
\bfm{a} = [a_1, a_2, \ldots, a_{n_{\text{out}}}],
$$
a \hl{max-pooling} layer with pooling window of size $r_p$ and stride $s_p$ produces the output
$$
a_i^{(\text{pool})} = \max_{j \in \mathcal{N}(i)} a_j, \quad i = 1, 2, \ldots, n_{\text{pool}},
$$
where $\mathcal{N}(i)$ denotes the $i^{\rm th}$ pooling window, given by
$$
\mathcal{N}(i) = \{(i-1)s_p + 1, (i-1)s_p + 2, \ldots, (i-1)s_p + r_p\},
$$
and
$$
n_{\text{pool}} = \Big\lfloor \frac{n_{\text{out}} - r_p}{s_p} \Big\rfloor + 1.
$$
If the maximum operation is replaced by the average, we obtain the \hl{average pooling} layer. Similarly, if the maximum operation is replaced by the minimum operation, then we obtain the \hl{min-pooling} layer.

Here, $\bfm{a}$ denotes the output feature map obtained from the preceding convolutional layer.  
The pooling operation acts on $\bfm{a}$ to produce a lower-dimensional representation 
$\bfm{a}^{(\text{pool})}$, which can be fed to the next convolutional or fully connected layer.
Schematically, we have
\[
\textbf{Input} 
\;\xrightarrow{\text{}}\;
\textbf{Convolution Layer} 
\;\xrightarrow{\text{}}\;
\textbf{Pooling}
\;\xrightarrow{\text{}}\;
\textbf{Next Layer}.
\]
  Thus, pooling serves two main purposes, namely,
\begin{itemize}
    \item It reduces the spatial resolution (or length) of the feature map, thereby decreasing the computational cost and the number of parameters in the subsequent fully connected layers. The pooling layer itself has no learnable parameters.
    \item It increases \hl{translational invariance}, meaning that small shifts or distortions in the input do not significantly change the pooled output.
\end{itemize}
\end{remark}
%%%%%%%%%%%%%%%%%%%%%%%%%%%%%%%%%%%%%%%%%
\begin{example}[1D Max Pooling]
Let the convolution output be
\[
\bfm{a} = [a_1, a_2, a_3, a_4, a_5, a_6],
\]
and consider max pooling with pooling size $r_p=2$ and stride $s_p=2$.  
Then,
\[
\begin{aligned}
a^{(\text{pool})}_1 &= \max(a_1, a_2),\\
a^{(\text{pool})}_2 &= \max(a_3, a_4),\\
a^{(\text{pool})}_3 &= \max(a_5, a_6).
\end{aligned}
\]
Hence, the pooled output is
\[
\bfm{a}^{(\text{pool})} = [a^{(\text{pool})}_1,\, a^{(\text{pool})}_2,\, a^{(\text{pool})}_3],
\]
which has reduced length compared to the original feature map.
\end{example}
%%%%%%%%%%%%%%%%%%%%%%%%%%%%%%%%%%%%%%%%%
Let us now illustrate translational invariance through pooling.
%%%%%%%%%%%%%%%%%%%%%%%%%%%%%%%%%%%%%%%%%
\begin{example}
Consider the convolutional output
$$
\bfm{a} = [5,\, 0,\, 4,\, 0].
$$
Applying max pooling with window size $2$ and stride $2$ gives
$$
\bfm{a}^{(\text{pool})} = [\max(5,0),\, \max(4,0)] = [5,\, 4].
$$
If the input is shifted by one position to the right, then, by translation equivariance, the feature map is also shifted by one position, producing
$$
\bfm{a}' = [0,\, 5,\, 0,\, 4],
$$
then the pooled output remains
$$
\bfm{a}^{\prime(\text{pool})} = [\max(0,5),\, \max(0,4)] = [5,\, 4].
$$
Hence, pooling increases \hl{translational invariance}, meaning that small shifts
in the input signal do not significantly affect the pooled output. A shift by two positions gives $\bfm{a}'' = [0,\, 0,\, 5,\, 0]$ and the pooled output $[0,\, 5]$. Hence, the invariance holds only for small shifts.

\end{example}
%%%%%%%%%%%%%%%%%%%%%%%%%%%%%%%%%%%%%%%%%
\begin{note}
In this example, the pooled output does not change, because each maximum stays inside its pooling window after the shift. If a maximum moves into another window, the pooled output changes. For instance, $\bfm{a} = [5,\, 0,\, 4,\, 100]$ gives $[5,\, 100]$, whereas its shift $[0,\, 5,\, 0,\, 4]$ gives $[5,\, 4]$. Hence, max pooling gives only an approximate invariance to small shifts.
\end{note}
%%%%%%%%%%%%%%%%%%%%%%%%%%%%%%%%%%%%%%%%%
\subsubsection{Multilayer Convolutional Networks}
%%%%%%%%%%%%%%%%%%%%%%%%%%%%%%%%%%%%%%%%%
A convolutional neural network is obtained by composing several convolutional layers, in the same way as a feedforward network is obtained by composing dense layers in \chap{feed.forward.architecture.ch}. For simplicity, we consider layers with one kernel. The superscript $(l)$ denotes the layer index, as in \chap{network.gradient.evaluation.ch}.
%%%%%%%%%%%%%%%%%%%%%%%%%%%%%%%%%%%%%%%%%
\begin{definition}[Convolutional Neural Network]
Let $L_c\in\N$. For $l = 1, 2, \ldots, L_c$, let $\bfm{w}^{(l)}\in\mathbb{R}^{r^{(l)}}$ be a kernel, $b^{(l)}\in\mathbb{R}$ a bias and $\actf^{(l)}$ an activation function. Let the stride and the padding of layer $l$ be given. A \defn{convolutional neural network} with $L_c$ convolutional layers computes
$$
\bfm{a}^{(0)} = \bfm{x},
\qquad
\bfm{a}^{(l)} = \mathcal{P}^{(l)}\Big(\vactf^{(l)}\big(\bfm{w}^{(l)} \star \bfm{a}^{(l-1)} + b^{(l)}\big)\Big),
\quad l = 1, 2, \ldots, L_c,
$$
where
\begin{itemize}
\item $\bfm{w}^{(l)} \star \bfm{a}^{(l-1)}$ is computed with the stride and the padding of layer $l$,
\item the bias $b^{(l)}$ is added to every entry,
\item $\vactf^{(l)}$ applies $\actf^{(l)}$ to every entry, and
\item $\mathcal{P}^{(l)}$ is a pooling operation or the identity map.
\end{itemize}
The vector $\bfm{a}^{(L_c)}$ is the input of a feedforward network, which gives the output
$$
\bfm{y} = \ann{L,\hat{n}}\big(\bfm{a}^{(L_c)};\bfm{\Theta}\big).
$$
The trainable parameters are the kernels $\bfm{w}^{(l)}$ and the biases $b^{(l)}$, together with the parameters $\bfm{\Theta}$ of the feedforward network.
\end{definition}
%%%%%%%%%%%%%%%%%%%%%%%%%%%%%%%%%%%%%%%%%
\begin{note}
The \hl{depth} of the network is $L_c + L$. It counts all layers with weights, including the output layer, as in \chap{feed.forward.architecture.ch}. Pooling layers are not counted, since they have no parameters.
\end{note}
%%%%%%%%%%%%%%%%%%%%%%%%%%%%%%%%%%%%%%%%%
The length of $\bfm{a}^{(l)}$ is obtained from the length of $\bfm{a}^{(l-1)}$ by the output length formula in the remark on padding, followed by the formula for pooling. The network is trained by backpropagation (\chap{network.gradient.evaluation.ch}) and the optimization methods of \chap{optimization.algorithms.ch}. Since a kernel entry is shared by all neurons of a layer, its gradient is the sum of the contributions of all these neurons.
%%%%%%%%%%%%%%%%%%%%%%%%%%%%%%%%%%%%%%%%%
The following example illustrates a one-dimensional convolutional neural network with $L_c = 1$ that predicts a time series, namely the price of a stock. Here, we choose the parameters by hand just for illustration purposes. In practice, the parameters are obtained by training, using the methods of \chap{network.gradient.evaluation.ch} and \chap{optimization.algorithms.ch}.
%%%%%%%%%%%%%%%%%%%%%%%%%%%%%%%%%%%%%%%%%
\begin{example}[A One-Dimensional CNN for Stock Price Prediction]
The closing prices of a stock on seven consecutive days are
$$
\bfm{x} = [100,\, 102,\, 101,\, 104,\, 103,\, 107,\, 106].
$$
We predict the change in the closing price on the next day with the following network.

\textbf{Convolutional layer.} Take the difference kernel $\bfm{w}^{(1)}=[-1,\,1]$ of the edge-detection example, with bias $b^{(1)}=0$, stride $1$, no padding and activation $\actf^{(1)} = \ReLU$. Then
$$
a_i = \ReLU\big(x_{i+1}-x_i\big),\qquad i=1,2,\ldots,6.
$$
The daily changes are $[2,\,-1,\,3,\,-1,\,4,\,-1]$. Hence, the feature map is
$$
\bfm{a} = [2,\, 0,\, 3,\, 0,\, 4,\, 0].
$$
Each entry is the rise of the price on one day. A fall gives the value zero.

\textbf{Pooling layer.} The pooling operation $\mathcal{P}^{(1)}$ is max pooling with window size $r_p=2$ and stride $s_p=2$. It gives
$$
\bfm{a}^{(1)} = [\max(2,0),\, \max(3,0),\, \max(4,0)] = [2,\, 3,\, 4].
$$
Each entry is the largest rise in a window of two days.

\textbf{Output layer.} The feedforward network consists of the output layer only, so that $L = 1$, $n_0 = 3$, $n_1 = 1$ and $\hat{n} = 0$. With weights $\bfm{v}=(0.1,\, 0.2,\, 0.3)$, bias $c=0$ and the identity activation, it gives the prediction
$$
\hat{y} = \bfm{v}\cdot\bfm{a}^{(1)} + c = 0.2 + 0.6 + 1.2 = 2.0.
$$
Thus, the network (with manually chosen parameters) predicts that the price rises by $2$, that is, from $106$ to $108$. The larger weights on the later entries give more importance to recent days.
\end{example}
%%%%%%%%%%%%%%%%%%%%%%%%%%%%%%%%%%%%%%%%%
\begin{remark}[Training and Practical Models]
In the above example, the kernel and the weights are chosen by hand. In practice, the parameters $\bfm{w}^{(1)}$, $b^{(1)}$, $\bfm{v}$ and $c$ are learned from historical data. Each training example $(\bfm{x}_i, y_i)$ consists of a window $\bfm{x}_i$ of past prices together with the observed change $y_i$ on the following day. The parameters minimize the mean squared error
$$
\mse\big(\bfm{w}^{(1)}, b^{(1)}, \bfm{v}, c\big) = \frac{1}{N_{\rm train}}\sum_{i=1}^{N_{\rm train}} \big(\hat{y}(\bfm{x}_i) - y_i\big)^2,
$$
where $\hat{y}(\bfm{x}_i)$ is the prediction of the network for the input $\bfm{x}_i$.
Practical models use log returns instead of prices. They also use several kernels, several input channels (such as price and volume) and several layers. 
\end{remark}
%%%%%%%%%%%%%%%%%%%%%%%%%%%%%%%%%%%%%%%%%

%%%%%%%%%%%%%%%%%%%%%%%%%%%%%%%%%%%%%%%%%
\subsection{Two-Dimensional Convolutional Layer}
%%%%%%%%%%%%%%%%%%%%%%%%%%%%%%%%%%%%%%%%%
We now extend the one-dimensional convolutional layer to the two-dimensional case,
which is the core operation in convolutional neural networks (CNNs) for image processing.
%%%%%%%%%%%%%%%%%%%%%%%%%%%%%%%%%%%%%%%%%
\subsubsection{Convolution Operator}
%%%%%%%%%%%%%%%%%%%%%%%%%%%%%%%%%%%%%%%%%
In this subsection, we briefly recall the basic definitions concerning convolution. We start with convolution of two functions from $\R^2$ to $\R$.
%%%%%%%%%%%%%%%%%%%%%%%%%%%%%%%%%%%%%%%%%
\begin{definition}[Convolution of Functions]

Let $f, g : \mathbb{R}^2 \to \mathbb{R}$. The \hl{convolution} of $f$ and $g$, denoted by $(f * g)$, is defined by
$$
(f * g)(x, y)
= \int_{\mathbb{R}^2} f(\xi, \eta)\, g(x - \xi,\, y - \eta)\, d\xi\, d\eta,
$$
whenever the integral exists. 
\end{definition}
%%%%%%%%%%%%%%%%%%%%%%%%%%%%%%%%%%%%%%%%%
The continuous definition of the convolution operator can be restricted to the discrete case, leading to the definition of convolution between two matrices.
%%%%%%%%%%%%%%%%%%%%%%%%%%%%%%%%%%%%%%%%%
\begin{definition}[Convolution of Matrices]

Let $A = [a_{i,j}] \in \mathbb{R}^{m \times n}$ and $W = [w_{p,q}] \in \mathbb{R}^{r_h \times r_w}$. 
The \hl{discrete convolution} of $A$ and $W$, denoted by $(W * A)$, is defined by
$$
(W * A)_{i,j}
= \sum_{p=1}^{r_h} \sum_{q=1}^{r_w}
w_{p,q}\, a_{\,i - p + 1,\, j - q + 1},
$$
where $a_{i - p + 1,\, j - q + 1} = 0$ whenever the indices fall outside the range of $A$.
\end{definition}
%%%%%%%%%%%%%%%%%%%%%%%%%%%%%%%%%%%%%%%%%
CNNs typically use the \hl{cross-correlation} operation given by
$$
(W \star A)_{i,j}
= \sum_{p=1}^{r_h} \sum_{q=1}^{r_w}
w_{p,q}\, a_{\,i + p - 1,\, j + q - 1}.
$$
%%%%%%%%%%%%%%%%%%%%%%%%%%%%%%%%%%%%%%%%%
\begin{note}
In the machine learning literature, the above cross-correlation itself is referred to as convolution. We follow this convention, and we write $W * A$ for $W \star A$ from now on.
\end{note}
%%%%%%%%%%%%%%%%%%%%%%%%%%%%%%%%%%%%%%%%%
\begin{remark}
The matrix $A$ can be viewed as an infinite matrix 
$$
A = [a_{i,j}]_{i,j=-\infty}^\infty
$$
with finite support, which means that $a_{i,j}=0$ whenever $i \notin \{1, \ldots, m\}$ or $j \notin \{1, \ldots, n\}$. In this case, we say that $A$ has an $m\times n$ support.
\end{remark}
%%%%%%%%%%%%%%%%%%%%%%%%%%%%%%%%%%%%%%%%%
\begin{example}[Moving average]

Given a two dimensional discrete signal $A$, the $2\times 2$ supported kernel that produces a simple moving average is given by
$$
W = \left[\begin{array}{cccccc}
	\tfrac{1}{4}&\tfrac{1}{4}\\
	\tfrac{1}{4}&\tfrac{1}{4}\\
	\end{array}\right],
$$
where $w_{1,1}=w_{1,2}=w_{2,1}=w_{2,2}=1/4$, and the matrix is extended by zeros in all four directions. A kernel whose entries are all equal and sum to one is called a \hl{mean filter} (or \hl{box filter}).  Then the cross-correlation leads to the simple moving average
$$
(W * A)_{i,j} = \sum_{p=1}^{r_h} \sum_{q=1}^{r_w} 
w_{p,q}\, a_{\,i + p - 1,\, j + q - 1} = \frac{1}{4}\big(a_{i,j} + a_{i+1,j} + a_{i,j+1} + a_{i+1,j+1} \big).
$$
This operation replaces each pixel value by the average of the entries of the $2\times 2$ block of $A$ given by
$$
B_{i,j} =
\begin{bmatrix}
a_{i,j} & a_{i,j+1}\\
a_{i+1,j} & a_{i+1,j+1}
\end{bmatrix}.
$$
This produces a \hl{blurred} or \hl{smoothed} version of the image.

Averaging replaces each pixel by the mean of its local neighborhood, which suppresses rapid spatial variations (high-frequency components) while preserving slowly varying content (low-frequency components). Consequently, sharp transitions (edges) and noise are attenuated, yielding a \hl{blurred} (smoothed) image. This relates it to the concept of a \hl{low-pass filter}.
\end{example}
%%%%%%%%%%%%%%%%%%%%%%%%%%%%%%%%%%%%%%%%%
\begin{note}
In the theoretical definition, convolution is expressed as an infinite sum 
ranging from $-\infty$ to $\infty$, assuming that the sequences are defined for all integer indices.  
In practical settings, however, the input and kernel have finite support, and all values outside their valid ranges are treated as zero.  
For computational convenience, we will omit the explicit zero extensions and write the summations only over the appropriate finite index ranges.
\end{note}
%%%%%%%%%%%%%%%%%%%%%%%%%%%%%%%%%%%%%%%%%
\begin{problem}[Blurring Effect of a Mean Filter]
Consider a $3\times3$ mean filter (or box filter) kernel
\[
W = \frac{1}{9}
\begin{bmatrix}
1 & 1 & 1\\
1 & 1 & 1\\
1 & 1 & 1
\end{bmatrix}.
\]
Show that the operation produces a blurred image (equivalently, show that the kernel $W$ is a low-pass filter).
\end{problem}
%%%%%%%%%%%%%%%%%%%%%%%%%%%%%%%%%%%%%%%%%
\begin{hint}
\begin{itemize}
\item Write down the expression for $(W * X)_{i,j}$, where $X = [x_{p,q}]$ is a two-dimensional input signal (or image) and show that each output pixel $a_{i,j}$ is the average of the $3\times3$ neighborhood of $X$ whose top-left entry is $(i,j)$.
\item Explain why this operation produces a \hl{blurred image} by giving an argument similar to the one given in the above example.
\end{itemize}
\end{hint}
%%%%%%%%%%%%%%%%%%%%%%%%%%%%%%%%%%%%%%%%%
%%%%%%%%%%%%%%%%%%%%%%%%%%%%%%%%%%%%%%%%%
\begin{example}[Gaussian Blurring]

A smoother and more natural way to blur an image is to use a \hl{Gaussian blur kernel},
whose entries are proportional to a discrete approximation of a two-dimensional Gaussian function.  
A commonly used $3\times3$ Gaussian blur kernel is
$$
W_{\text{G}} = \frac{1}{16}
\begin{bmatrix}
1 & 2 & 1\\
2 & 4 & 2\\
1 & 2 & 1
\end{bmatrix}.
$$
Convolving an image $X$ with $W_{\text{G}}$ produces the output
$$
(W_{\text{G}} * X)_{i,j}
= \sum_{p=1}^{3}\sum_{q=1}^{3}
w_{p,q}\, x_{i+p-1,\,j+q-1},
$$
where $w_{p,q}$ are the entries of $W_{\text{G}}$.

\medskip
The weights are largest at the center and decrease symmetrically outward,
so this filter performs a \hl{weighted average} that preserves overall image brightness
while reducing noise and small-scale variations.  
Unlike the simple mean filter, the Gaussian kernel avoids sharp artifacts and
produces a more visually natural \hl{blurred} image.
\end{example}
%%%%%%%%%%%%%%%%%%%%%%%%%%%%%%%%%%%%%%%%%
The mean filter smooths (blurs) an image by averaging nearby pixels, whereas
an \hl{edge detection} filter highlights regions of sharp intensity change as illustrated in the following example. 
%%%%%%%%%%%%%%%%%%%%%%%%%%%%%%%%%%%%%%%%%
\begin{example}[Edge Detection via Convolution]

A simple \hl{vertical edge detection kernel} is
$$
W_v =
\begin{bmatrix}
-1 & 0 & 1\\
-1 & 0 & 1\\
-1 & 0 & 1
\end{bmatrix},
$$
and a corresponding \hl{horizontal edge detection kernel} is
$$
W_h =
\begin{bmatrix}
-1 & -1 & -1\\
0 & 0 & 0\\
1 & 1 & 1
\end{bmatrix}.
$$
For an input image $X = [x_{i,j}]$, the feature maps are computed as
$$
A_v = W_v * X,
\qquad
A_h = W_h * X.
$$
Each of these responds strongly where there is a vertical or horizontal edge, respectively.

The matrices $W_v$ and $W_h$ are called \hl{Prewitt filters}. The \hl{Sobel filters} give the weight $2$ to the central row or column, as in the next example.
 
\medskip
The overall edge magnitude can be approximated by combining both directions:
$$
A_{\text{edge}} = \sqrt{A_v \odot A_v + A_h \odot A_h},
$$
where $\odot$ denotes the Hadamard product introduced in \chap{network.gradient.evaluation.ch}, and the square root is taken entrywise.
\end{example}
%%%%%%%%%%%%%%%%%%%%%%%%%%%%%%%%%%%%%%%%%
The following example contrasts Gaussian blurring with edge detection by the Sobel and Laplacian kernels.
%%%%%%%%%%%%%%%%%%%%%%%%%%%%%%%%%%%%%%%%%
\begin{example}[Blurring vs. Edge Detection]
The \hl{Gaussian blur} smooths an image by averaging pixel intensities with Gaussian weights.
For instance, using the $3\times3$ Gaussian kernel
$$
W_{\text{G}} = \frac{1}{16}
\begin{bmatrix}
1 & 2 & 1\\
2 & 4 & 2\\
1 & 2 & 1
\end{bmatrix},
$$
the output
\[
(W_{\text{G}} * X)_{i,j}
= \sum_{p=1}^{3}\sum_{q=1}^{3} w_{p,q}\,x_{i+p-1,\,j+q-1},
\]
where $w_{p,q}$ are the entries of $W_{\text{G}}$, replaces each pixel by a weighted local average.  
This operation suppresses high-frequency components, resulting in a \hl{blurred (low-pass)} image.

\medskip
In contrast, an \hl{edge detection} filter accentuates regions of sharp intensity change.
Two common examples are the \hl{Sobel} and \hl{Laplacian} kernels.

\begin{itemize}
\item \textbf{Sobel kernels:}
\[
W_x =
\begin{bmatrix}
-1 & 0 & 1\\
-2 & 0 & 2\\
-1 & 0 & 1
\end{bmatrix},
\qquad
W_y =
\begin{bmatrix}
-1 & -2 & -1\\
0 & 0 & 0\\
1 & 2 & 1
\end{bmatrix}.
\]
These measure horizontal and vertical intensity gradients, respectively. For an input image $X$, the feature maps are
\[
A_x = W_x * X,
\qquad
A_y = W_y * X,
\]
and the edge magnitude is
\[
A_{\text{edge}} = \sqrt{A_x \odot A_x + A_y \odot A_y},
\]
where the square root is taken entrywise.

\item \textbf{Laplacian kernel:}
\[
W_{\text{L}} =
\begin{bmatrix}
0 & -1 & 0\\
-1 & 4 & -1\\
0 & -1 & 0
\end{bmatrix}.
\]
This kernel approximates the negative of the Laplacian, $-\Delta$, of the image intensity. The Laplacian emphasizes regions where intensity changes rapidly in multiple directions,
producing strong responses along edges and contours.
\end{itemize}

\medskip
Thus, while the Gaussian blur removes detail by averaging neighboring pixels,
edge-detection filters highlight rapid local variations.
Together, these operations illustrate the fundamental duality between
\hl{low-pass filtering (smoothing)} and \hl{high-pass filtering (edge detection)} in image processing.
\end{example}
%%%%%%%%%%%%%%%%%%%%%%%%%%%%%%%%%%%%%%%%%

%%%%%%%%%%%%%%%%%%%%%%%%%%%%%%%%%%%%%%%%%
\subsubsection{2D Convolution Layer}
%%%%%%%%%%%%%%%%%%%%%%%%%%%%%%%%%%%%%%%%%
Let us restrict to grayscale images, where the input and the kernel are two-dimensional arrays (matrices).

Let the input to the layer be a two-dimensional array
$$
X = [x_{i,j}]_{i=1,j=1}^{n_{0h},n_{0w}} \in \mathbb{R}^{n_{0h}\times n_{0w}},
$$
where $n_{0h}$ and $n_{0w}$ denote the \hl{height} and \hl{width} of the input.

Let the convolutional kernel (or filter) be
$$
W = [w_{p,q}]_{p=1,q=1}^{r_h, r_w} \in \mathbb{R}^{r_h\times r_w},
$$
where $r_h$ and $r_w$ denote the kernel height and width, respectively.
Let $b \in \mathbb{R}$ be the bias term and $\actf$ be a nonlinear activation function.
%%%%%%%%%%%%%%%%%%%%%%%%%%%%%%%%%%%%%%%%%
\begin{definition}[Two-Dimensional Convolutional Layer]

For given stride values $s_h, s_w \ge 1$ in the vertical and horizontal directions, respectively,
the \hl{convolutional layer} produces the output
$$
a_{i,j} = \actf\!\Bigg(
\sum_{p=1}^{r_h}\sum_{q=1}^{r_w}
w_{p,q}\, x_{(i-1)s_h + p,\,(j-1)s_w + q} +\,b
\Bigg),
\quad
i=1,2,\ldots,n_{\text{out},h}, \;
j=1,2,\ldots,n_{\text{out},w},
$$
where
$$
n_{\text{out},h} = 
\Big\lfloor \frac{n_{0h} - r_h}{s_h} \Big\rfloor + 1,
\qquad
n_{\text{out},w} =
\Big\lfloor \frac{n_{0w} - r_w}{s_w} \Big\rfloor + 1,
$$
are the output dimensions. The pair $(s_h, s_w)$ is called the \hl{stride}.
\end{definition}
%%%%%%%%%%%%%%%%%%%%%%%%%%%%%%%%%%%%%%%%%

%%%%%%%%%%%%%%%%%%%%%%%%%%%%%%%%%%%%%%%%%
\begin{remark}[Padding]

If zero-padding of size $p_h$ and $p_w$ is applied to the height and width directions,
the effective input dimensions become
$$
\tilde{n}_{0h} = n_{0h} + 2p_h, \qquad \tilde{n}_{0w} = n_{0w} + 2p_w,
$$
and the output dimensions are given by
$$
n_{\text{out},h} = 
\Big\lfloor \frac{\tilde{n}_{0h} - r_h}{s_h} \Big\rfloor + 1,
\qquad
n_{\text{out},w} =
\Big\lfloor \frac{\tilde{n}_{0w} - r_w}{s_w} \Big\rfloor + 1.
$$
\end{remark}
%%%%%%%%%%%%%%%%%%%%%%%%%%%%%%%%%%%%%%%%%
\begin{example}[2D Convolution with $3\times3$ Kernel and Stride $(1,1)$]
Consider the input image and the kernel, respectively,
\[
X =
\begin{bmatrix}
x_{1,1} & x_{1,2} & x_{1,3} & x_{1,4} \\
x_{2,1} & x_{2,2} & x_{2,3} & x_{2,4} \\
x_{3,1} & x_{3,2} & x_{3,3} & x_{3,4} \\
x_{4,1} & x_{4,2} & x_{4,3} & x_{4,4}
\end{bmatrix},
\qquad
W =
\begin{bmatrix}
w_{1,1} & w_{1,2} & w_{1,3} \\
w_{2,1} & w_{2,2} & w_{2,3} \\
w_{3,1} & w_{3,2} & w_{3,3}
\end{bmatrix}.
\]
Assume that the kernel $W$ slides over $X$ with stride $(s_h, s_w) = (1,1)$ 
and without padding.  
Since $ n_{0h}= n_{0w} = 4$ and $r_h = r_w = 3$, the output dimensions are
\[
n_{\text{out},h} = n_{\text{out},w} = (4 - 3)/1 + 1 = 2.
\]
Thus, the output activation map $A = [a_{i,j}] \in \mathbb{R}^{2\times2}$ is given by
\[
a_{i,j} = \actf\!\Bigg(
\sum_{p=1}^{3}\sum_{q=1}^{3}
w_{p,q}\, x_{i+p-1,\,j+q-1} +\,b
\Bigg),
\qquad i,j = 1,2.
\]
\end{example}
%%%%%%%%%%%%%%%%%%%%%%%%%%%%%%%%%%%%%%%%%
\begin{problem}[Fixed Kernels in Convolutional Layers]
Consider a two-dimensional convolutional layer with the identity activation and zero bias,
$$
a_{i,j} = \sum_{p=1}^{r_h}\sum_{q=1}^{r_w} w_{p,q}\,x_{i+p-1,\,j+q-1}.
$$
Suppose the kernel $W$ is fixed as the $3\times3$ Gaussian kernel
$$
W_{\text{G}} = \frac{1}{16}
\begin{bmatrix}
1 & 2 & 1\\
2 & 4 & 2\\
1 & 2 & 1
\end{bmatrix}.
$$
Show that this layer performs Gaussian smoothing (blurring) of the input.
Explain what changes if the identity activation is replaced by \ReLU.
\end{problem}
%%%%%%%%%%%%%%%%%%%%%%%%%%%%%%%%%%%%%%%%%

%%%%%%%%%%%%%%%%%%%%%%%%%%%%%%%%%%%%%%%%%
\begin{problem}[Computing Output Feature Map Dimensions]
In each of the following problems, determine the output feature map size $(n_{\text{out},h}, n_{\text{out},w})$ of the two-dimensional convolutional layer defined above, with zero padding as in the remark on padding. The input size is $(n_{0h}, n_{0w})$, the kernel size is $(r_h, r_w)$, the stride is $(s_h, s_w)$, and the padding is $(p_h, p_w)$.

\begin{enumerate}
\item
An image of size $(32, 32)$ is convolved with a $3\times3$ Gaussian blur kernel
using stride $(1, 1)$ and no padding $(p_h = p_w = 0)$.

\item
The same image is convolved with a $5\times5$ kernel
using stride $(2, 2)$ and no padding $(p_h = p_w = 0)$.

\item
Now apply a $3\times3$ Laplacian edge-detection kernel
using stride $(1, 1)$ and padding $(p_h = p_w = 1)$ (“same” padding).
Show that the output has the same spatial dimensions as the input.

\item
Let an input of size $(64, 64)$ pass through two convolutional layers in sequence:
\begin{itemize}
\item Layer 1: kernel size $(3, 3)$, stride $(1, 1)$, padding $(1, 1)$,
\item Layer 2: kernel size $(5, 5)$, stride $(2, 2)$, padding $(2, 2)$.
\end{itemize}
Compute the final output feature-map size after both layers.
\end{enumerate}
\end{problem}
%%%%%%%%%%%%%%%%%%%%%%%%%%%%%%%%%%%%%%%%%

%%%%%%%%%%%%%%%%%%%%%%%%%%%%%%%%%%%%%%%%%
\begin{remark}[Pooling in Two Dimensions]

A two-dimensional \hl{max-pooling layer} with window size $r_{p,h}\times r_{p,w}$ and stride $(s_{p,h}, s_{p,w})$
acts on the feature map $A = [a_{i,j}]$ as
$$
a^{(\text{pool})}_{i,j}
= \max_{(k,l)\in\mathcal{N}_{i,j}} a_{k,l},
$$
where
$$
\mathcal{N}_{i,j} = \big\{\big((i-1)s_{p,h} + u,\, (j-1)s_{p,w} + v\big) : u = 1, 2, \ldots, r_{p,h},~ v = 1, 2, \ldots, r_{p,w}\big\}
$$
is the $(i,j)^{\rm th}$ pooling window.
This reduces the spatial resolution while improving translational invariance.

Similarly, \hl{min-pooling layer} and \hl{average-pooling layer} can also be defined.
\end{remark}
%%%%%%%%%%%%%%%%%%%%%%%%%%%%%%%%%%%%%%%%%
\begin{remark}[Colour Images and Multiple Channels]
So far, the input has been a grayscale image, which has one intensity value at each pixel. A colour image has three values at each pixel, one for each of the colours red, green and blue (RGB). Such an image is stored as an array
$$
X = [x_{i,j,c}] \in \mathbb{R}^{n_{0h}\times n_{0w}\times C},\qquad C = 3,
$$
where $c$ is the \hl{channel} index. For a fixed $c$, the matrix $[x_{i,j,c}]$ is called the $c^{\rm th}$ channel of $X$. A kernel for such an input has the same number of channels and is given by
$$
W = [w_{p,q,c}] \in \mathbb{R}^{r_h\times r_w\times C}.
$$
The output of the layer at position $(i,j)$ is
$$
a_{i,j} = \actf\Bigg(\sum_{c=1}^{C}\sum_{p=1}^{r_h}\sum_{q=1}^{r_w} w_{p,q,c}\, x_{(i-1)s_h + p,\,(j-1)s_w + q,\,c} + b\Bigg).
$$

The above layer has one kernel and produces one feature map. In practice, a layer has several kernels of the same size. Their number $K$ is chosen by the designer of the network, in the same way as the width of a layer in a feedforward network. We write the entries of the $k^{\rm th}$ kernel as $w_{p,q,c,k}$ and its bias as $b_k$. Each kernel is applied to the same input and gives its own feature map:
$$
a_{i,j,k} = \actf\Bigg(\sum_{c=1}^{C}\sum_{p=1}^{r_h}\sum_{q=1}^{r_w} w_{p,q,c,k}\, x_{(i-1)s_h + p,\,(j-1)s_w + q,\,c} + b_k\Bigg),\qquad k = 1, 2, \ldots, K.
$$
The $K$ feature maps together form an output with $K$ channels, which is the input of the next layer. The entries of all kernels are learned in training. Such a layer has $K(r_h r_w C + 1)$ parameters. This number does not depend on the size of the image.
\end{remark}
%%%%%%%%%%%%%%%%%%%%%%%%%%%%%%%%%%%%%%%%%
\begin{remark}[Multilayer Networks in Two Dimensions]
The convolutional neural network defined for one-dimensional inputs carries over to two dimensions. The feature maps $A^{(l)}$ are then matrices, with $A^{(0)} = X$ and
$$
A^{(l)} = \mathcal{P}^{(l)}\Big(\vactf^{(l)}\big(W^{(l)} * A^{(l-1)} + b^{(l)}\big)\Big),
\quad l = 1, 2, \ldots, L_c,
$$
where $W^{(l)}\in\mathbb{R}^{r^{(l)}_h\times r^{(l)}_w}$ is the kernel of layer $l$. Before the feedforward network, the matrix $A^{(L_c)}$ is \hl{flattened}, that is, its entries are listed row by row to form a vector. For inputs with several channels, each layer uses several kernels, as in the remark on colour images and multiple channels. Each $A^{(l)}$ is then an array with several channels.
\end{remark}
%%%%%%%%%%%%%%%%%%%%%%%%%%%%%%%%%%%%%%%%%
\begin{example}[A Small CNN for Digit Images]
Consider grayscale images of handwritten digits of size $28\times28$, which are to be classified into the $10$ classes $0, 1, \ldots, 9$. Take a network with two convolutional layers. Each layer has a $5\times5$ kernel, stride $1$, no padding and the activation \ReLU. Each layer is followed by max pooling with a $2\times2$ window and stride $2$. The sizes are as follows.
\begin{itemize}
\item Layer 1: the convolution gives $(28-5)+1 = 24$, and pooling gives $\lfloor (24-2)/2\rfloor + 1 = 12$. Hence, $A^{(1)}\in\mathbb{R}^{12\times12}$.
\item Layer 2: the convolution gives $(12-5)+1 = 8$, and pooling gives $\lfloor (8-2)/2\rfloor + 1 = 4$. Hence, $A^{(2)}\in\mathbb{R}^{4\times4}$.
\item Flattening $A^{(2)}$ gives a vector in $\mathbb{R}^{16}$. An output layer with $10$ neurons and the softmax function $\texttt{sm}$ gives the class probabilities.
\end{itemize}
The number of parameters is $25+1 = 26$ for each convolutional layer and $16\times10 + 10 = 170$ for the output layer, that is, $222$ in total. For comparison, a single dense layer from the $784$ pixels to the $10$ outputs has $784\times10 + 10 = 7850$ parameters.
\end{example}
%%%%%%%%%%%%%%%%%%%%%%%%%%%%%%%%%%%%%%%%%
\begin{remark}
A network with one kernel per layer is too small for practical use. Practical networks use several kernels in each layer, as described in the remark on colour images and multiple channels. For example, the network LeNet-5  has the same structure as the above example, with $6$ kernels in its first convolutional layer and $16$ kernels in its second. We omit these details for this course.
\end{remark}
%%%%%%%%%%%%%%%%%%%%%%%%%%%%%%%%%%%%%%%%%
\subsection{A Real Architecture: VGG-16}\label{vgg16.subs}
%%%%%%%%%%%%%%%%%%%%%%%%%%%%%%%%%%%%%%%%%
We conclude the discussion of CNNs with a network that is used in practice. VGG-16 was introduced by Simonyan and Zisserman in 2014. It classifies colour images of size $224\times224\times3$ into the $1000$ classes of the ImageNet data set, such as animals, vehicles and everyday objects. Its architecture is shown in \figu{vgg16.fig}.
%%%%%%%%%%%%%%%%%%%%%%%%%%%%%%%%%%%%%%%%%
\begin{figure}[ht]
\begin{center}
\input{Figures/Ch08VGG16CNN.png}
\end{center}
\caption{The architecture of VGG-16. The height of each box shows the spatial size of the feature map, and its width shows the number of channels.}\label{vgg16.fig}
\end{figure}
%%%%%%%%%%%%%%%%%%%%%%%%%%%%%%%%%%%%%%%%%
In the notation of \sect{convolution.neural.network.sec}, VGG-16 is a convolutional neural network with the following structure.
\begin{enumerate}
\item \textbf{Convolutional layers.} There are $13$ convolutional layers, arranged in five blocks of $2$, $2$, $3$, $3$ and $3$ layers. Every layer uses kernels of spatial size $3\times3$, stride $1$, padding $1$ and the activation \ReLU. Hence, every convolutional layer keeps the spatial size of its input.
\item \textbf{Channels.} Each layer has several kernels, as in the remark on colour images and multiple channels. The number of kernels is $64$ in the first block, and it doubles in each of the next blocks up to $512$. For example, each kernel of the first layer has size $3\times3\times3$, and each kernel of the second layer has size $3\times3\times64$.
\item \textbf{Pooling.} Each block ends with max pooling with a $2\times2$ window and stride $2$. The pooling operation of every other layer is the identity map. The spatial size is halved by each pooling, from $224$ to $112$, $56$, $28$, $14$ and finally $7$.
\item \textbf{Feedforward network.} The output of size $7\times7\times512$ is flattened into a vector of length $25\,088$. A feedforward network with two hidden layers of width $4096$ and the activation \ReLU, and an output layer with $1000$ neurons, maps this vector to the class scores. The softmax function $\texttt{sm}$ turns these scores into the probabilities of the $1000$ classes.
\end{enumerate}
The depth of VGG-16 is $13 + 3 = 16$, which gives the network its name. The pooling layers are not counted, since they have no parameters. The network has about $138$ million parameters. Most of them, about $103$ million, belong to the first dense layer. All parameters were learned on about $1.2$ million labelled ImageNet images, using backpropagation and gradient-based optimization as in \chap{network.gradient.evaluation.ch} and \chap{optimization.algorithms.ch}.
%%%%%%%%%%%%%%%%%%%%%%%%%%%%%%%%%%%%%%%%%
\begin{autref}
VGG-16: K. Simonyan and A. Zisserman, Very deep convolutional networks for large-scale image recognition, arXiv:1409.1556 (2014). Add a bib entry if needed. The cat photograph in \figu{vgg16.fig} is \texttt{Figures/imagescat.jpg}; check its source and licence before distributing the notes.
\end{autref}
%%%%%%%%%%%%%%%%%%%%%%%%%%%%%%%%%%%%%%%%%
\begin{code}
The Colab notebook
\href{https://drive.google.com/file/d/11ddy1PA-9L2SbUFrRdsuoSZqNwa0HePd/view?usp=sharing}{\texttt{\color{thc}{ImageClassificationVGG16}}}
applies the trained VGG-16 to an image of your choice. It
\begin{enumerate}
\item loads VGG-16 with the parameters learned on ImageNet,
\item preprocesses the uploaded image into an input of size $224\times224\times3$,
\item prints the five most probable classes and their probabilities, and
\item plots the $64$ feature maps of the first two convolutional layers.
\end{enumerate}
Run the notebook with a few images of your own. Locate in the code the preprocessing, the output of the feedforward part and the feature maps. Also try an image whose main object is not one of the ImageNet classes, for example a person, and explain the result.
\end{code}
%%%%%%%%%%%%%%%%%%%%%%%%%%%%%%%%%%%%%%%%%
